%% Section 3: Neural Recording Modalities
%% Complete content integrated from original paper with NO EM DASHES and active voice

\section{Neural Recording Modalities}
\label{sec:modalities}

\subsection{The Architecture of Signal Access}

Every brain-to-speech system begins at the same decision point: how to listen to the brain. The choice of neural recording modality shapes everything that follows. It determines the spatial granularity of the signal, the temporal fidelity with which speech dynamics can be captured, the risk the participant must accept, and ultimately the ceiling on what any downstream AI architecture can achieve. No amount of algorithmic sophistication can recover information that was never present in the recording. This section characterizes the six primary neural recording modalities and three hybrid or emerging approaches represented across the 71 studies included in this review, grounding each in the neurophysiology of speech production and perception that underpins the decoding task. Figure~\ref{fig:modality_matrix} summarizes the spatial-temporal resolution tradeoffs across all modalities; Figure~\ref{fig:electrode_schematics} illustrates the physical placement and electrode configurations for each approach.

Speech generation does not occur in a single cortical region but in a distributed, hierarchically organized network. The dual-stream model \cite{Hickok2007,Poeppel2020} describes a dorsal pathway connecting auditory cortex to sensorimotor regions, responsible for the sensorimotor transformations of speech production. A ventral pathway connects auditory cortex to the frontal lobe via the temporal lobe, supporting speech comprehension and phonological working memory. Both pathways contribute signals of interest to brain-to-speech decoding, and the modality employed determines which parts of this network are accessible \cite{Herff2016,Rabbani2019}. The choice of modality is therefore not merely an engineering decision. It is a decision about which aspects of the speech-language system the AI decoder will have the opportunity to learn from.

\subsection{Invasive Modalities}

Invasive neural recording requires surgical access to the brain. In exchange for that irreversible commitment, invasive modalities provide signal quality that non-invasive approaches cannot match. Higher spatial resolution, broader frequency bandwidth, substantially higher signal-to-noise ratio, and direct access to the cortical regions most tightly coupled to speech motor execution characterize these approaches. The landmark clinical results of the past five years, communication rates approaching natural speech and word error rates below 10 percent, have been achieved exclusively with invasive modalities. Whether that clinical performance can ultimately be matched without the surgical requirement remains the defining open question of the field.

\subsubsection{Electrocorticography (ECoG)}

Electrocorticography involves placing an electrode array directly on the cortical surface following a craniotomy, recording the local field potentials generated by the aggregate synaptic and neuronal activity of underlying cortical columns. Unlike scalp EEG, ECoG electrodes are positioned beneath the skull and dura, eliminating the spatial low-pass filtering that severely degrades surface recordings. Standard clinical ECoG grids typically have 4 mm inter-electrode spacing and 2.3 mm electrode diameters. High-density research arrays achieve 1 mm or submillimeter spacing over selected cortical regions. This spatial resolution enables the differentiation of signals from adjacent cortical areas, a capability critical for decoding speech, where the phonotopic and articulatory maps of the sensorimotor cortex have spatial structure on the scale of millimeters.

The temporal bandwidth of ECoG is equally valuable. High-gamma activity (70-150 Hz) proves particularly informative for speech decoding \cite{Chartier2018,Ramsey2018}, reflecting the firing of neuronal ensembles that modulate reliably with phonological and articulatory events. The envelope of the high-gamma band correlates well with speech production kinematics. Its rapid temporal dynamics allow the system to track the fine structure of running speech at rates approaching natural articulation, approximately 4-6 syllables per second in fluent English. Lower-frequency oscillations (alpha at 8-13 Hz and beta at 13-30 Hz) also carry speech-relevant information, particularly for imagined and covert speech where the high-gamma signal is weaker.

The research history of ECoG-based speech decoding spans the full timeline of this review. The Brain-to-Text study \cite{Herff2015} established the proof of principle, achieving approximately 25 percent WER on continuous spoken phrases from seven participants. This notable result relied on classical hidden Markov model speech recognition. Subsequent work \cite{Chartier2018} demonstrated that articulatory kinematic trajectories can be encoded in sensorimotor cortex with Pearson correlations reaching $r=0.87$, providing the neurophysiological basis for the articulatory intermediate representation approach that Anumanchipalli and colleagues \cite{Anumanchipalli2019} would later exploit in their landmark Nature paper, which produced the first fluent and intelligible ECoG-synthesized speech. The machine translation analogy \cite{Makin2020} then introduced the encoder-decoder framework that enabled direct cortical activity-to-text decoding at approximately 3 percent WER on a constrained corpus.

ECoG has also proved particularly valuable for studying imagined and covert speech. Martin and colleagues \cite{Martin2016} demonstrated above-chance classification of imagined word pairs from ECoG. Proix and colleagues \cite{Proix2022} showed that imagined speech can be decoded from low- and cross-frequency ECoG features, revealing a neural oscillatory basis for inner speech distinct from the high-gamma activity dominant in overt speech. The high-density ECoG approach \cite{Ramsey2018}, employing 256-channel grids at 1 mm spacing, demonstrated phoneme-level decoding accuracy of 75.5 percent across five participants, establishing that increased spatial sampling of the sensorimotor cortex yields commensurate improvements in phonological decoding. The largest ECoG cohort assembled to date \cite{Chen2024Neural}, trained across 48 participants, achieving Pearson correlations of $r=0.804$ for acoustic reconstruction, demonstrated for the first time that ECoG decoding scales to population-level data.

The principal limitation of standard ECoG is its requirement for craniotomy. Most research ECoG studies are conducted in epilepsy patients undergoing clinical seizure monitoring, implanted for clinical rather than research purposes. This population's neural characteristics may not fully generalize to the ALS and paralysis populations that speech neuroprostheses are designed to serve. Metzger and colleagues \cite{Metzger2023}, using high-density ECoG in a participant with anarthria following brainstem stroke, demonstrated that this gap can be bridged. However, the surgical requirement remains an irreducible entry barrier for the majority of potential beneficiaries.

\subsubsection{Intracortical Microelectrode Arrays}

Where ECoG samples the aggregate activity of cortical columns, intracortical microelectrode arrays penetrate the cortical grey matter itself, recording the extracellular electrical signatures of individual neurons or small neuronal clusters within millimeters of the electrode tip \cite{Wilson2020}. The Utah array, a 10 by 10 grid of 1.5 mm silicon shanks with 400 micrometer inter-electrode spacing, is the most widely deployed implant in the BrainGate program. It is the device underlying the most clinically significant speech decoding results published to date. The richness of the single-unit and multi-unit activity available from intracortical recordings provides a fundamentally higher-bandwidth communication channel between brain and decoder than any surface modality can offer.

The clinical results from the intracortical modality represent the current frontier of the field. Moses and colleagues \cite{Moses2021} demonstrated the first real-time sentence decoding in a person with anarthria at 15 words per minute. Willett and colleagues \cite{Willett2023} extended this to 62 words per minute at 9 percent WER for a 50-word vocabulary, a communication rate approaching that of impaired natural speech. Card and colleagues \cite{Card2024} achieved 2.5 percent WER sustained over eight months with a 30-minute daily calibration protocol. Wairagkar and colleagues \cite{Wairagkar2025} demonstrated an instantaneous voice synthesis neuroprosthesis capable of reproducing paralinguistic features including intonation and prosody in real time. Collectively, these results establish intracortical recording as the currently dominant modality for clinical speech neuroprosthesis performance.

Several properties of intracortical recordings beyond raw signal quality contribute to this dominance. The threshold crossing rates of motor cortical neurons are tightly coupled to attempted articulatory movements, providing a high-bandwidth, rapidly modulating signal even in individuals who can no longer produce any overt speech output. Degenhart and colleagues \cite{Degenhart2020} demonstrated that low-dimensional manifold alignment between sessions preserves approximately 95 percent of decoding performance without re-collecting labeled data, addressing the practical challenge of daily signal drift that all chronically implanted systems face. Wandelt and colleagues \cite{Wandelt2024} showed that individual neurons in the supramarginal gyrus, an area involved in phonological working memory, encode internal speech representations even in patients who have been unable to speak for years, extending the reach of intracortical decoding beyond primary motor cortex to the broader speech-language network.

The primary limitation of intracortical arrays is biological. Microelectrodes induce a chronic inflammatory response, glial scarring, that progressively attenuates signal quality over months to years. Signal loss and electrode failure are significant concerns for long-term clinical deployment \cite{Vansteensel2023}. The surgical risk of intracortical implantation, while manageable, is higher than that of epidural ECoG. The FDA-regulated approval pathway for such devices is correspondingly complex, and the number of participants who have undergone research intracortical implantation remains small.

\subsubsection{Stereotactic EEG and Depth Electrodes (sEEG)}

Stereotactic EEG occupies a compelling middle ground in the invasiveness landscape. Thin multi-contact depth electrodes, typically 0.8-1.0 mm in diameter with contact separations of 3.5 mm, are implanted via small burr holes without craniotomy, traversing the brain to sample cortical and subcortical structures along their trajectory. The surgical footprint is substantially smaller than that of ECoG, and the ability to sample deep structures (the insula, hippocampus, and basal ganglia) that are inaccessible to surface electrodes makes sEEG uniquely informative about the broader cortical-subcortical networks involved in speech production and comprehension.

The minimally invasive nature of sEEG has begun to attract growing research attention for speech BCIs. Angrick and colleagues \cite{Angrick2021} demonstrated real-time synthesis of imagined speech from frontal cortex sEEG signals using a CNN-neural vocoder pipeline. Kohler and colleagues \cite{Kohler2022} showed that encoder-decoder architectures applied to sEEG depth electrodes can achieve acoustic reconstruction quality comparable to ECoG. Most significantly, Singh and colleagues \cite{Singh2025} assembled the largest sEEG speech decoding cohort reported in the literature, pooling distributed recordings from 25 participants to demonstrate that a group decoder trained across participants outperforms any individual participant decoder, establishing cross-participant generalization as achievable for sEEG, not merely for surface modalities. The sEEG modality represents 4 of the 32 invasive studies in this review, but its clinical accessibility and unique access to deep structures suggest its representation in the literature will grow substantially.

\subsection{Non-Invasive Modalities}

Non-invasive modalities offer a fundamentally different proposition: access to neural signals without any surgical procedure, with correspondingly lower signal quality and a different set of constraints on what can be decoded and how. For the hundreds of millions of individuals with speech or language impairments who are ineligible for surgical intervention, or who would not choose it, non-invasive BCI represents the only viable neural pathway to communication assistance. The scientific and engineering challenges are formidable, but the clinical imperative is equally compelling.

\subsubsection{Scalp Electroencephalography (EEG)}

Scalp EEG measures the sum of postsynaptic currents from millions of synchronously active cortical neurons via electrodes placed on the scalp surface, typically in a 64- or 128-channel configuration following the standardized 10-20 system. The skull, cerebrospinal fluid, and scalp tissue act collectively as a spatial low-pass filter, attenuating high-frequency components, blurring the spatial specificity of cortical sources, and mixing signals from distinct generators into composite scalp recordings that are difficult to unmix. The signal-to-noise ratio is substantially lower than ECoG across all frequency bands \cite{Schirrmeister2017}, and the high-gamma range, so informative for invasive speech decoding, is effectively inaccessible due to muscle artifact contamination above approximately 40 Hz.

Despite these constraints, EEG has attracted considerable research attention for brain-to-speech decoding, driven by its accessibility, low cost, and the absence of surgical risk. Early EEG speech decoding work focused on closed-set imagined speech classification, identifying which of a small set of words or phonemes a participant was imagining, typically using support vector machines or linear discriminant analysis applied to spectral features. The deep learning era brought convolutional and recurrent architectures that improved classification accuracies. Saha and colleagues \cite{Saha2019} demonstrated subject-independent hierarchical deep learning achieving approximately 70-75 percent accuracy across ten participants. Agarwal and Kumar \cite{Agarwal2022} demonstrated 73.56 percent accuracy for five-class imagined speech with a subject-independent LSTM across 13 participants.

A qualitative shift occurred with the introduction of transformer architectures and language model integration into the EEG decoding pipeline. Wang and Ji \cite{Wang2022} demonstrated the first open-vocabulary EEG-to-text system using a BART language model, reporting above-chance sentence reconstruction from EEG during reading. Duan and colleagues \cite{Duan2023} extended this with the DeWave architecture, using discrete EEG wave encodings that do not require explicit time-marker alignment, a practical advance that simplified training and improved BLEU and ROUGE metrics over the prior state of the art. Yet it is important to note that EEG open-vocabulary WER in these paradigms remains approximately 35-45 percent, substantially higher than what ECoG or intracortical systems achieve, reflecting the fundamental signal quality ceiling of scalp recording. The majority of EEG studies in this review use closed-set vocabularies, reflecting the modality's current limitations for unconstrained communication.

The EEG modality accounts for 20 of the 71 included studies, the second-largest group after invasive ECoG and intracortical approaches, spanning imagined speech classification \cite{Panachakel2021,LopezBernal2022,Nieto2022}, perceived speech decoding, and emerging open-vocabulary generation paradigms. Its accessibility makes it the most practical non-invasive option for clinical populations who cannot access more specialized recording facilities.

\subsubsection{Magnetoencephalography (MEG)}

Magnetoencephalography measures the tiny magnetic fields, on the order of femtotesla $(10^{-15}$ T), generated by the same neuronal current flows that produce the EEG signal, using superconducting quantum interference device (SQUID) sensors housed in a liquid helium-cooled cryostat. Because magnetic fields pass through the skull and scalp without the attenuation or spatial distortion that affects electrical potentials, MEG offers substantially better spatial source localization than EEG, on the order of 2-3 mm for superficial cortical sources with femtotesla sensitivity \cite{Baillet2001}, whilst retaining the millisecond temporal resolution that is the shared strength of electrophysiological modalities. MEG is selectively sensitive to tangentially oriented current sources, providing complementary cortical coverage to EEG, which is more sensitive to radial sources.

For brain-to-speech decoding, MEG has proven particularly powerful when combined with large participant samples and machine learning approaches that exploit its improved spatial specificity. The landmark study \cite{Defossez2023} applied contrastive self-supervised learning to MEG and EEG recordings from 175 participants across four public datasets \cite{Gwilliams2023}, achieving 41 percent top 1 classification accuracy for perceived speech segment identification. This scale of evaluation was only possible because multiple public MEG datasets existed to enable it, a circumstance that reflects the MEG research community's stronger tradition of open data sharing compared to the invasive literature. Kwon and colleagues \cite{Kwon2024} demonstrated the first intelligible speech synthesis directly from MEG signals using a transformer and BigVGAN neural vocoder, achieving Mel-spectrogram Pearson correlations of approximately $r=0.95$. Most recently, d'Ascoli and colleagues \cite{dAscoli2025} trained a single CNN-transformer decoder across 723 participants from seven datasets covering English, French, and Dutch, demonstrating that MEG word decoding performance scales with the number of training participants, a finding with important implications for the scalability of non-invasive neural speech decoding.

MEG's principal limitation is practical rather than physiological: the requirement for a magnetically shielded room and the prohibitive cost of the cryogenic infrastructure render it unsuitable for clinical deployment outside specialist research centers \cite{Boto2018}. The introduction of optically pumped magnetometer (OPM) arrays, which operate at room temperature and can be worn as a helmet, may eventually address this barrier. However, OPM-based speech decoding studies did not appear in the included literature during the review period.

\subsubsection{Functional MRI (fMRI)}

Functional MRI measures the blood-oxygen-level-dependent (BOLD) signal, the haemodynamic response to neural activity arising from the coupling between neuronal firing and local increases in cerebral blood flow and oxygenation. With a temporal resolution of 1-2 seconds (set by the repetition time of the acquisition sequence) and a spatial resolution of 1-3 mm voxels, fMRI occupies the opposite corner of the resolution trade-off space from EEG: superb spatial specificity across the entire brain simultaneously, but temporal resolution that is orders of magnitude coarser than the neural events of interest. The haemodynamic response peaks approximately 4-5 seconds after the neural event it reflects \cite{Logothetis2010,Hemodynamic2007}, a fundamental physiological latency that no signal processing approach can fully overcome.

For speech decoding, this temporal limitation means that fMRI cannot track the millisecond-to-millisecond dynamics of speech articulation. What it can do, uniquely among non-invasive modalities, is reconstruct the semantic and linguistic content of perceived and imagined speech at the level of meaning rather than sound, an application that has yielded some of the most notable results in the non-invasive brain-to-speech literature. Tang and colleagues \cite{Tang2023Semantic} demonstrated continuous semantic language reconstruction from fMRI signals using a GPT-based alignment approach, achieving 100 percent story identification accuracy and substantial above-chance recovery of narrative content. The breadth and specificity of the semantic maps revealed by fMRI \cite{Huth2016}, provide the neural substrate for this capability: the meaning of language is genuinely represented in the BOLD signal, even if the acoustic form is not.

Recent work has extended fMRI decoding to cross-participant paradigms. Tang and Huth \cite{Tang2025Semantic} demonstrated cross-participant semantic decoding from fMRI signals across different stimulus modalities (story-based and movie-based presentation), establishing that shared semantic representations generalize both across individuals and across stimulus types. The fMRI modality contributes 6 of the 14 MEG/fMRI studies in this review, with the remainder using MEG; the proportion of fMRI studies has increased markedly from 2022 onwards, coinciding with the availability of large open fMRI language datasets and the maturation of LLM-based decoding pipelines that are particularly well-suited to the semantic character of fMRI signals.

\subsection{Hybrid and Emerging Modalities}

Beyond the canonical six modalities, the included literature encompasses three hybrid and emerging approaches that address specific limitations of their unimodal counterparts by combining complementary signal sources, increasing spatial sampling density, or exploiting peripheral articulatory signals that are not derived from the brain itself but are tightly coupled to speech production intent.

\subsubsection{EEG plus fNIRS Fusion}

Functional near-infrared spectroscopy (fNIRS) measures cortical haemodynamics via the differential absorption of near-infrared light at oxyhaemoglobin and deoxyhaemoglobin chromophores, providing a slow (1-2 second temporal resolution) but portable and wearable measure of cortical blood oxygenation. When combined with EEG in a hybrid recording system, the two modalities provide complementary information: EEG captures rapid neural oscillations (millisecond resolution, poor spatial specificity) whilst fNIRS contributes slower haemodynamic context (low temporal resolution, better cortical spatial specificity than EEG alone). Cooney and colleagues \cite{Cooney2022} demonstrated that a bimodal CNN trained on EEG+fNIRS data outperforms unimodal EEG decoding for both overt and imagined speech across 19 participants. Hons and colleagues \cite{Hons2026} further extended EEG+fNIRS hybrid decoding to phoneme classification from imagined and perceived speech, confirming the complementarity of the combined signal. The practical advantage of EEG+fNIRS over EEG alone comes at modest additional cost in hardware and setup complexity, making it a promising direction for wearable speech BCIs.

\subsubsection{Surface Electromyography (sEMG) for Silent Speech}

Surface electromyography is not, strictly speaking, a neural modality. It records the electrical activity of facial and laryngeal muscles rather than the brain itself. However, it occupies a legitimate place in the brain-to-speech BCI landscape because it captures the articulatory intent of speech at the peripheral motor execution stage: the muscle patterns associated with speaking a word silently (mouthed or sub-vocal speech) are highly consistent and informative, even without overt vocalization. High-density sEMG (HD-sEMG) arrays with 120 or more surface electrodes across facial and submental muscle groups provide particularly rich spatial sampling of articulatory dynamics.

The HD-sEMG results in this review are among the highest-performing in the non-invasive literature. Song and colleagues \cite{Song2023Silent} achieved 97.17 percent phrase accuracy and 3.11 percent character error rate for 33 Chinese phrases using a spatiotemporal CNN-RNN architecture on HD-sEMG, the best CER reported for any silent speech approach in this review. Zhu and colleagues \cite{Zhu2021} demonstrated 87-90 percent digit classification accuracy for both English and Chinese using an optimized 120-electrode HD-sEMG configuration. Li and colleagues \cite{Li2022} extended HD-sEMG to Mandarin tonal speech synthesis using an encoder-decoder architecture. These results suggest that peripheral articulatory signals, whilst technically not 'brain' signals in the neural sense, represent a highly accessible and high-performing avenue for silent speech communication that deserves consideration alongside brain-based approaches.

\subsubsection{Micro-ECoG (microECoG)}

Micro-ECoG refers to high-density ECoG arrays with sub-millimeter electrode spacing and electrode diameters on the order of 100-200 micrometers, fabricated using flexible polymer substrates rather than conventional rigid silicon or polyimide grids. The spatial sampling advantage of micro-ECoG over standard ECoG is directly analogous to the advantage of intracortical recording over ECoG: finer spatial resolution reveals higher-bandwidth information about the cortical speech motor map that is averaged away at coarser scales.

Duraivel and colleagues \cite{Duraivel2023} provided the most direct quantitative comparison of micro-ECoG and standard macro-ECoG for speech decoding, demonstrating a 35 percent improvement in phoneme classification accuracy from micro-ECoG versus macro-ECoG in the same four participants, attributable to the 57-fold increase in effective spatial sampling. This result carries important implications: it suggests that a significant portion of the performance gap between ECoG and intracortical systems may be recoverable by increasing ECoG spatial density, potentially offering a less biologically disruptive path to higher-performance speech decoding. Micro-ECoG devices face their own engineering challenges (fabrication complexity, the need for high channel count recording systems, and limited coverage area per device) but represent an active direction in next-generation neuroprosthesis hardware.

\subsection{Comparative Summary}

Table 1 in the main text and Figures~\ref{fig:modality_matrix} and \ref{fig:electrode_schematics} summarize the technical, operational, and performance characteristics of all modality groups represented across the 71 included studies. Several patterns are immediately apparent. First, the performance gradient from intracortical recording through ECoG and sEEG to non-invasive modalities is consistent and substantial. The best WER achieved by any non-invasive EEG system (approximately 35-45 percent for open-vocabulary tasks) remains roughly 10-15x larger than the best intracortical result (2.5 percent, Card and colleagues). When compared to typical invasive performance (9-15\% WER), non-invasive EEG shows a 3-5x performance gap, reflecting fundamental signal quality constraints. Second, the study distribution is markedly skewed towards invasive modalities in the early years of the review period (2015-2019), with non-invasive and hybrid modalities gaining representation from 2022 onwards, reflecting the enabling role of transformer architectures and self-supervised learning in making non-invasive decoding more competitive. Third, non-invasive modalities dominate the cross-participant generalization evidence (detailed in Section 6), suggesting that the larger participant samples accessible without surgical barriers enable the multi-subject training that generalization requires.

No modality in this review is universally superior. The choice is always a function of the clinical context, the patient population, and the specific communication goal. A completely locked-in ALS patient who requires immediate communication restoration and is willing to accept surgical risk stands to benefit most from intracortical implantation, where the performance evidence is most compelling. A patient with progressive aphasia who retains some residual communication capacity and prefers to avoid surgery may derive practical benefit from a well-designed EEG+fNIRS or EEG imagined speech system at lower performance levels. The systematic evidence synthesized here is intended to ground those clinical decisions in the actual empirical record, rather than in the most recently published headline result.

